\chapter*{Resumen}
\addcontentsline{toc}{chapter}{Resumen}
\markboth{Resumen}{Resumen}

Este Trabajo de Fin de Grado aborda el desarrollo de un banco de pruebas basado en simulación física y fotorrealista para evaluar algoritmos de localización y mapeo simultáneo visual (SLAM visual) junto con navegación autónoma en entornos industriales y espacios abiertos. El entorno desacopla la simulación sensorial y cinemática, ejecutada en NVIDIA Isaac Sim~6.2 con trazado de rayos en tiempo real, de la pila de percepción y control desplegada sobre ROS~2 Humble en contenedores Docker independientes conectados mediante la pila de red del sistema anfitrión.

El sistema se orienta a contrastar dos aproximaciones representativas de SLAM estéreo bajo trayectorias reproducibles y con referencia exacta de posición (\textit{ground truth}): por un lado, \texttt{ORB-SLAM3} como referencia clásica ejecutada sobre CPU; por otro, \texttt{NVIDIA Isaac cuVSLAM} como alternativa acelerada en GPU mediante CUDA. La evaluación cuantitativa de las trayectorias estimadas se plantea con la librería \texttt{evo} a partir de las métricas de error absoluto de pose (ATE) y error relativo de pose (RPE), integrando posteriormente el stack \texttt{Nav2} para cerrar el lazo de navegación autónoma y preparar la transferencia hacia plataformas reales del laboratorio de la UPV (TurtleBot3 en fase planar y vehículo aéreo en fase de seis grados de libertad).

En el punto actual del desarrollo se ha verificado la capacidad de cómputo de la estación de trabajo mediante pruebas preliminares de renderizado estéreo a \num{1280}\(\times\)\num{720}~píxeles y se ha completado la migración de \texttt{tfg\_ros2\_core} a \texttt{ament\_cmake} (C++17) en ROS~2 Humble, compilada sin errores ni advertencias mediante \texttt{colcon build} en \qty{6.69}{\second} (\qty{6.82}{\second} en el resumen global) con los ejecutables \texttt{stereo\_sync\_monitor} y \texttt{tum\_trajectory\_logger}. Esta base deja preparados la propagación de \texttt{use\_sim\_time:=true}, el lazo de teleoperación sobre \texttt{/cmd\_vel}, la monitorización estéreo, el registro de trayectorias en formato TUM, el perfil de RViz2 y el guion de captura en \texttt{rosbag2}, quedando pendientes las pruebas en vivo tras conectar el puente de Isaac Sim.

\positinconcluso{Redactar las versiones en inglés (\textit{Abstract}) y en valenciano (\textit{Resum}) si la normativa de la titulación en la UPV las exige, e incorporar al finalizar el proyecto las cifras experimentales definitivas de ATE, RPE, FPS reales y consumo de recursos.}
